% The experimental report (experiments.md, "The report").
% Build:  cd report && latexmk -lualatex -outdir=build report.tex
% Tables: sbt "campaign/run report" writes tables/*.tex and tables/prompt.txt from
%         the attempt records; never edit those by hand. Shawn's inspection lives
%         in inspection.json, which the report reads and never writes.
\documentclass[11pt]{article}
\usepackage{fontspec}
\setmonofont{DejaVu Sans Mono}[Scale=MatchLowercase] % covers the contract's σ′, ⊥, ↦
\usepackage[margin=2.2cm]{geometry}
\usepackage{booktabs,graphicx,pdflscape,fvextra}
\usepackage[hidelinks]{hyperref}

% A generated file, or a note saying how to make it. Each inclusion is logged, so
% a build can be checked for showing real tables rather than placeholders.
\newcommand{\generated}[2]{\IfFileExists{tables/#1}{\typeout{pag-report: included #1}#2}%
  {\emph{Not generated yet: run \texttt{sbt "campaign/run report"}.}}}

\title{One-shot abstract domains from small models\\\large E1, rung 0: the Qwen3.5 size ladder}
\author{Shawn Meier}
\date{\today}

\begin{document}
\maketitle

\section{Question}
% TODO: at what model size can a small open-weight model write a sound, useful
% backward abstract domain in one shot, from the contract alone?

\section{Setup}
% TODO: the models (Qwen3.5 0.8B to 27B, own Q8 GGUFs, llama.cpp -c 65536 -np 1);
% the prompt (rung 0 of the information ladder: the task in general terms and the
% contract; Appendix~\ref{app:prompt}); the smoke corpus, never shown to a model;
% one shot, no feedback.

\begin{table}[ht]
  \centering
  \caption{The settings each campaign ran with, from its records. Sampling as
    the Qwen3.5-27B model card recommends for coding in thinking mode;
    ``default'' is the server's own.}
  \label{tab:settings}
  \resizebox{\linewidth}{!}{\generated{settings.tex}{\input{tables/settings.tex}}}
\end{table}

\begin{table}[ht]
  \centering
  \caption{The smoke corpus: one \texttt{reach} target per probe, with its true
    answer. Reachable targets are run with inputs that reach them, so a domain
    refuting one is caught unsound.}
  \label{tab:corpus}
  \small\generated{corpus.tex}{\input{tables/corpus.tex}}
\end{table}

\section{Results}

\begin{landscape}
\begin{table}[p]
  \centering
  \caption{One row per attempt. Cells: R refuted, A alarm, $\times$ refuted a
    reachable target (unsound), I did not converge, E domain failure, H hung.
    The last three columns are Shawn's inspection.}
  \label{tab:attempts}
  \resizebox{\linewidth}{!}{\generated{table1.tex}{\input{tables/table1.tex}}}
\end{table}
\end{landscape}

\begin{table}[ht]
  \centering
  \caption{Per model. Mechanical: evaluated, not unsound, proves at least one
    target. Inspection: accepted out of those inspected. Median proved: among
    the mechanically acceptable. The rest: where attempts stopped.}
  \label{tab:models}
  \resizebox{\linewidth}{!}{\generated{table2.tex}{\input{tables/table2.tex}}}
\end{table}

\section{Observations}
% TODO: written by hand, from the inspection notes.

\appendix
\section{The prompt}\label{app:prompt}
Exactly as sent, from the attempt records.
\generated{prompt.txt}{\VerbatimInput[fontsize=\scriptsize,breaklines,breakanywhere]{tables/prompt.txt}}

\end{document}
